%% kapa: exam=1 round=35 year=2024 period=1 form=A subject=economics subject_ko=경제학원론 seq=37 number=77
%% source: Q-Net 감정평가사 기출문제 (cst00309) artlSeq=5213229 file=raw/1008/5213229/1교시 문제지 원본.pdf
%% conversion: kordoc 4.24.51 -> tools/gen1.py (auto); review=proofread-sonnet; flags=none
\begin{question}{77}
현재 한국과 미국의 햄버거 가격이 각각 5,000원, 5달러인 경우, 이에 관한 설명으로 옳은 것을 모두 고른 것은? (단, 햄버거를 대표상품으로 한다.)
\begin{framedtext}
ㄱ. 현재 구매력평가 환율은 1,000(원/달러)이다. \par
ㄴ. 변동환율제도하에서 현재 환율이 1,100(원/달러)이다. 장기적으로 구매력 평가설이 성립하고 미국의 햄버거 가격과 환율이 변하지 않는다면, 장기적으로 한국의 햄버거 가격은 상승한다. \par
ㄷ. 변동환율제도하에서 현재 환율이 1,100(원/달러)이다. 장기적으로 구매력 평가설이 성립하고 한국과 미국의 햄버거 가격이 변하지 않는다면, 장기적으로 환율은 상승한다. \par
\end{framedtext}
\begin{choices}
  \item ㄱ
  \item ㄷ
  \item ㄱ, ㄴ
  \item ㄴ, ㄷ
  \item ㄱ, ㄴ, ㄷ
\end{choices}
\end{question}
\answer{3}
